% !TeX root = ../thuthesis-example.tex

\chapter{引言}

% 本章仅包含原文的第一节：引言

\label{sec:intro}
随着《中华人民共和国个人信息保护法》《中华人民共和国数据安全法》等法律法规的相继出台，以及《中共中央 国务院关于构建数据基础制度更好发挥数据要素作用的意见》明确提出将数据要素作为推动数字经济发展的关键资源，数据资产的管理与价值计量已成为企业数字化转型的核心议题。然而，在推荐系统中，以深度兴趣网络\cite{zhou2018deep}为代表的深度模型将高维行为数据压缩为排序或匹配预测，使系统效益改进与底层数据单元之间的因果关系呈现明显“黑箱”特征。传统资产估值方法（成本法、收益法、市场法）因数据资产的非实体性、价值易变性与强情境依赖而难以直接适配\footnote{参见《加快构建数据资产估值体系 赋能数字经济发展》，\url{https://www.deo.org.cn/article/hygd_sjpg/detail143.html}}。如何在合规框架下对数据要素进行可计算的价值度量，已成为数据市场与数据治理领域的重要研究议题\cite{xu2023data}。

在上述背景下，本文将问题先收敛到\textbf{归因层}，即先回答“哪些数据单元在模型泛化中产生了可重复、可解释的边际影响”。这样处理主要基于三点考虑。首先在方法层面，样本梯度与验证梯度可形成可计算的局部影响代理，便于在统一训练轨迹下进行重复检验。在工程层面，归因指标可以直接复用既有检查点与验证流程，减少为单一估值任务进行大规模重训练的成本。在治理层面，归因结论更容易附带统计边界与不确定性说明，也更符合“先证据、后决策”的审慎披露要求。基于这一思路，本文将“归因可信度”作为连接模型效果与数据价值讨论的起点。

学术界已提出多类数据价值评估方法，但各有其局限。基于合作博弈的数据Shapley值\cite{ghorbani2019data}具有严格的公平性公理保证，但其精确计算需枚举所有$2^N$个数据子集，已被证明为$\#P$-hard问题\cite{andos2012computation}，在现代深度推荐模型中完全不可行。影响函数\cite{influence}通过局部二阶近似将"删除样本"操作转化为连续可微的权重扰动问题，无需重训练即可估计样本影响，但其依赖精确Hessian逆矩阵的计算——对于参数量达$10^7$量级的推荐模型，存储完整Hessian需$O(d^2)$空间，求逆需$O(d^3)$时间，工程上难以落地。TracIn\cite{TracIn}通过沿训练轨迹累积梯度内积估计样本影响，避免了Hessian求逆，但需保存完整训练轨迹或多个检查点，计算开销仍然较高（本文实验中耗时约27.9秒/次），且未能区分不同训练阶段的差异化贡献。上述方法的共同局限在于：它们或将训练过程视为静态黑箱，或在计算效率与估值精度之间存在难以调和的矛盾，难以满足工业级推荐系统对"实时、可解释、可部署"数据价值监控的需求。



